\documentclass[11pt]{article}

\usepackage[margin=1in]{geometry}
\usepackage{amsmath,amssymb,booktabs,longtable,enumitem,hyperref}
\usepackage[T1]{fontenc}

\title{Morphological Profiling to Gene Programs and Gene Sets}
\author{}
\date{}

\begin{document}
\maketitle

\section*{Overview}
This note defines a morphology-based extractor for converting Cell Painting or other imaging-derived morphological profiles into gene-level weights and gene sets. The central use case is a \emph{query profile} measured in a user experiment, which is compared against a \emph{reference library} of profiles from perturbations with semantic handles (compounds, CRISPR knockouts, ORF overexpression). The reference library provides a route from profile similarity to gene hypotheses, and finally to ranked gene sets.

The intended implementation follows the same design principles as other extractors in the repository:
\begin{itemize}[leftmargin=2em]
\item a clear separation between raw assay processing and extractor logic,
\item conservative defaults that produce interpretable outputs,
\item optional resource bundles for reusable references,
\item and explicit metadata, warnings, and GMT output.
\end{itemize}

\section{Biological rationale}
Morphological profiling captures downstream consequences of perturbing cellular systems. A query profile can therefore be interpreted by asking which \emph{known perturbations} in a reference library induce similar morphologies. If the similar reference perturbations are genetic (CRISPR or ORF), the mapping to genes is direct. If they are compounds, the mapping is mediated by a compound-to-target table.

This yields a natural posterior-like gene scoring model:
\begin{enumerate}[leftmargin=2em]
\item compare a query morphology profile to reference perturbation profiles,
\item convert similarity to a perturbation evidence weight,
\item route perturbation evidence to genes through direct or target-based mappings,
\item aggregate evidence across perturbations and modalities,
\item convert the resulting gene weights to one or more gene sets.
\end{enumerate}

In practice, Cell Painting references such as JUMP are especially attractive because they contain large collections of compounds and genetic perturbations with strong metadata, including matched compound/target and gene perturbation panels. Public JUMP resources include a pilot set of 300+ compounds and 160+ genes in A549 and U2OS, a very large production set with over 116,000 compounds and 15,000+ genes in U2OS, and a curated JUMP-Target metadata resource linking compounds, ORFs, and CRISPR reagents to target genes. These resources are publicly documented by the JUMP consortium and the Cell Painting Gallery. \cite{jump_main,jump_gallery,jump_target,jump_pilot,jump_morphmap}

\section{Inputs and notation}
Let
\begin{itemize}[leftmargin=2em]
\item $f \in \{1,\dots,F\}$ index morphology features,
\item $q \in \{1,\dots,Q\}$ index query profiles or query groups,
\item $r \in \{1,\dots,R\}$ index reference perturbation profiles,
\item $g \in \{1,\dots,G\}$ index genes.
\end{itemize}

The query profile matrix is $X^{(Q)} \in \mathbb{R}^{Q \times F}$.
The reference profile matrix is $X^{(R)} \in \mathbb{R}^{R \times F}$.
The extractor assumes these profiles are already computed and aggregated at the appropriate level (for example, well-level normalized and feature-selected profiles from a CellProfiler or pycytominer workflow).

Reference metadata provides:
\begin{itemize}[leftmargin=2em]
\item perturbation type $t_r \in \{\mathrm{compound}, \mathrm{crispr}, \mathrm{orf}\}$,
\item a perturbation identifier $p_r$,
\item optional QC weight $c_r \in [0,1]$,
\item optional cell line, timepoint, and batch labels.
\end{itemize}

Compound target metadata provides a sparse mapping $A \in \mathbb{R}_{\ge 0}^{R \times G}$ where $A_{rg}$ is the weight by which compound perturbation $r$ implicates gene $g$. For genetic perturbations, a direct one-hot mapping $B \in \mathbb{R}_{\ge 0}^{R \times G}$ is used:
\[
B_{rg} =
\begin{cases}
1 & \text{if perturbation } r \text{ directly perturbs gene } g, \\
0 & \text{otherwise.}
\end{cases}
\]

\section{Feature alignment and standardization}
\subsection{Feature intersection}
Let $\mathcal{F}$ be the intersection of features present in the query and reference profiles. Only these features are used. The extractor should warn if the fraction of reference features recovered in the query is low.

\subsection{Reference-based scaling}
If the reference bundle provides feature means and scales (or robust medians and MADs), define standardized query and reference values:
\[
\tilde{X}^{(Q)}_{qf} = \frac{X^{(Q)}_{qf} - \mu_f}{\sigma_f + \epsilon}, \qquad
\tilde{X}^{(R)}_{rf} = \frac{X^{(R)}_{rf} - \mu_f}{\sigma_f + \epsilon}.
\]
If no explicit statistics are supplied, the implementation may assume the profiles are already normalized in a common feature space and leave them unchanged.

\section{Query aggregation}
If multiple query rows correspond to the same biological unit (for example, replicate wells or a user-supplied group), aggregate them before similarity calculation. A robust default is the componentwise median:
\[
\bar{X}^{(Q)}_{uf} = \mathrm{median}\{ \tilde{X}^{(Q)}_{qf} : q \in \mathcal{Q}(u) \},
\]
where $\mathcal{Q}(u)$ is the set of query profiles assigned to unit $u$.
This produces one extracted program per query unit $u$.

\section{Similarity model}
\subsection{Base similarity}
The default similarity between query unit $u$ and reference perturbation $r$ is cosine similarity:
\[
\mathrm{sim}(u,r) = \frac{\bar{X}^{(Q)}_u \cdot \tilde{X}^{(R)}_r}{\|\bar{X}^{(Q)}_u\|_2 \, \|\tilde{X}^{(R)}_r\|_2 + \epsilon}.
\]
An optional alternative is Pearson correlation. The implementation should support both.

\subsection{Polarity}
Morphological similarity can be interpreted in at least two ways:
\begin{itemize}[leftmargin=2em]
\item \textbf{similar} (default): retrieve perturbations with similar morphology,
\item \textbf{opposite}: retrieve perturbations with anti-correlated morphology,
\item \textbf{both}: emit both directions as separate programs.
\end{itemize}

Define positive and negative evidence:
\[
s^+_{ur} = \max(\mathrm{sim}(u,r), 0)^\gamma, \qquad
s^-_{ur} = \max(-\mathrm{sim}(u,r), 0)^\gamma,
\]
where $\gamma \ge 1$ is a similarity sharpness parameter.

\subsection{Perturbation QC weights}
If a perturbation-level QC weight $c_r$ is provided (for example, replicate consistency or profile activity), define:
\[
w^{+}_{ur} = c_r s^+_{ur}, \qquad
w^{-}_{ur} = c_r s^-_{ur}.
\]
If no QC weights are supplied, set $c_r = 1$.

\section{Routing perturbation evidence to genes}
\subsection{Compound perturbations}
For compounds, evidence is distributed across targets according to the compound target map:
\[
C_{rg} = A_{rg}.
\]
If only an unweighted target list is available, use uniform weights within each compound:
\[
\sum_g A_{rg} = 1 \quad \text{for each compound } r \text{ with at least one target.}
\]

\subsection{Genetic perturbations}
For genetic perturbations, evidence maps directly to the perturbed gene:
\[
G_{rg} = B_{rg}.
\]

\subsection{Modality-specific gene scores}
For query unit $u$, define the positive similarity gene evidence from compounds and genetics separately:
\[
E^{+,\mathrm{compound}}_{ug} = \sum_{r : t_r = \mathrm{compound}} w^+_{ur} C_{rg},
\]
\[
E^{+,\mathrm{genetic}}_{ug} = \sum_{r : t_r \in \{\mathrm{crispr}, \mathrm{orf}\}} w^+_{ur} G_{rg}.
\]

To prevent one modality from dominating simply because it has more perturbations, normalize each by its total evidence mass:
\[
\hat{E}^{+,\mathrm{compound}}_{ug} = \frac{E^{+,\mathrm{compound}}_{ug}}{\sum_h E^{+,\mathrm{compound}}_{uh} + \epsilon},
\]
\[
\hat{E}^{+,\mathrm{genetic}}_{ug} = \frac{E^{+,\mathrm{genetic}}_{ug}}{\sum_h E^{+,\mathrm{genetic}}_{uh} + \epsilon}.
\]

Then combine modalities with weights $\lambda_c$ and $\lambda_g$:
\begin{equation}
S^{+}_{ug} = \lambda_c \hat{E}^{+,\mathrm{compound}}_{ug} + \lambda_g \hat{E}^{+,\mathrm{genetic}}_{ug},
\label{eq:morph_positive_score}
\end{equation}
where by default $\lambda_c = \lambda_g = 1/2$ if both are available, and $\lambda=1$ for the available modality if only one exists.

An analogous score $S^-_{ug}$ is constructed from negative similarity evidence $w^-_{ur}$.

\section{Gene program extraction}
\subsection{Program families}
Each query unit $u$ yields one or more gene programs:
\begin{itemize}[leftmargin=2em]
\item \texttt{similar}: based on $S^+_{ug}$,
\item \texttt{opposite}: based on $S^-_{ug}$ if requested.
\end{itemize}
The default is to emit only the similar program.

\subsection{Selection}
Given a nonnegative gene score vector $S_{ug}$ for a chosen polarity, select genes by:
\begin{itemize}[leftmargin=2em]
\item \textbf{top-k} (default): choose the $K$ highest-scoring genes,
\item \textbf{quantile}: choose genes above a quantile threshold,
\item \textbf{threshold}: choose genes above a minimum score.
\end{itemize}
As in other extractors, the implementation should enforce default size bounds of 100--500 genes for GMT emission and warn when fewer genes are available.

\subsection{Neighbor restriction for specificity}
In practice, morphology retrieval can become too diffuse if all weakly similar references are allowed to contribute. A robust default is therefore to retain only the top $K_{\mathrm{nbr}}$ reference neighbors per query and polarity after thresholding and hubness-aware evidence reweighting:
\[
\mathcal{N}_u^+ = \operatorname{TopK}\{r : \mathrm{sim}(u,r) \ge \tau\},
\]
and analogously for the opposite-polarity case using $-\mathrm{sim}(u,r)$.
Only retained neighbors contribute to $E_{ug}$. This improves specificity by preventing a long tail of weak matches from dominating the routed gene evidence. In the practical implementation, hubness-adjusted evidence is computed first and the retained neighborhood is then truncated to a narrow local panel.

\subsection{Same-modality-first retrieval}
If the query metadata identifies the perturbation modality (for example compound versus ORF), same-modality neighbors can be prioritized before cross-modality fusion. Cross-modality neighbors are still useful, but they should contribute most strongly when they reinforce the same local target signal as the same-modality neighborhood rather than dominating it outright.

\subsection{Adaptive neighborhood}
In practice, $K_{\mathrm{nbr}}$ is better interpreted as an upper bound than as a fixed neighborhood size. If evidence drops sharply after a small coherent local panel has been assembled, the extractor may stop early rather than adding a long tail of weak or isolated neighbors. This improves specificity and makes the effective retained-neighbor count query-dependent.

\subsection{Within-set normalization}
For the selected set $\mathcal{S}_u$, define within-set weights:
\[
W_{ug} =
\begin{cases}
\displaystyle \frac{S_{ug}}{\sum_{h \in \mathcal{S}_u} S_{uh} + \epsilon} & g \in \mathcal{S}_u, \\
0 & \text{otherwise.}
\end{cases}
\]
This yields a probability-like distribution over selected genes, suitable for weighted enrichment tools.

\section{Reference bundles}
A morphology extractor is most useful when paired with a reusable reference bundle. The recommended v1 bundle pattern is:
\begin{itemize}[leftmargin=2em]
\item \texttt{reference\_profiles.parquet}: consensus perturbation profiles (one row per perturbation/context),
\item \texttt{reference\_metadata.tsv.gz}: perturbation type, cell line, timepoint, QC, control labels,
\item \texttt{compound\_targets.tsv.gz}: compound-to-gene target map with weights,
\item \texttt{feature\_schema.tsv.gz}: ordered feature list,
\item \texttt{feature\_stats.tsv.gz}: per-feature center and scale,
\item \texttt{bundle.json}: version, provenance, source URLs, checksums.
\end{itemize}

A practical first bundle should be derived from the JUMP pilot data in a single cell line and timepoint (for example, U2OS 48h or A549 48h), because those datasets contain matched compounds, CRISPR knockouts, and ORF overexpression perturbations with strong metadata and manageable scale compared to the full JUMP production release. \cite{jump_gallery,jump_pilot,jump_target}

The practical default should prefer coherent same-timepoint bundles. If a public context does not contain all modalities at one timepoint, the workflow should keep a coherent same-timepoint reference from the available modalities and warn explicitly about missing modalities. Mixed-timepoint bundles should require explicit user intent.

\subsection{Reference hubness penalty}
Some reference perturbations are generic morphology hubs: they are moderately similar to many unrelated queries and can dominate routed gene evidence. Let $h_r$ denote a precomputed hubness score for reference perturbation $r$, estimated from within-bundle reference-reference similarities. A default penalty rescales evidence before routing:
\[
\tilde{w}_{ur} = \frac{w_{ur}}{\epsilon + h_r}.
\]
The implementation may also use a rank-based penalty. In practice this is often more stable than a raw inverse-linear rescaling when hub scores are compressed. The intent is the same: reduce the influence of generic hub-like references and increase specificity of routed gene evidence.

\subsection{Gene recurrence penalty}
Some genes recur across many unrelated morphology perturbations because they are common library targets or generic downstream responders. A weak default penalty can be defined from reference-level recurrence:
\[
\omega_g = \frac{\log((N+1)/(d_g+1))}{\log(N+1)},
\]
where $d_g$ is the number of reference perturbations that route to gene $g$ and $N$ is the total number of mapped references. The final routed gene score can then be rescaled as
\[
\tilde{S}_{ug} = \omega_g S_{ug},
\]
which reduces recurrent generic genes without imposing a hard blacklist.

\section{QC and failure modes}
Important warnings to implement:
\begin{itemize}[leftmargin=2em]
\item Low feature overlap between query and reference.
\item Query profiles dominated by metadata columns or non-numeric fields.
\item Reference control profiles accidentally included as retrieval candidates.
\item Low perturbation QC or very low number of usable reference matches.
\item Gene families dominating selected sets (for example, a broad receptor family), reported as warnings but not filtered by default.
\item Low retrieval confidence, for example when few positive neighbors are retained or the negative-similarity fraction is high.
\item Low specificity confidence, for example when retained neighbors are geometrically coherent but routed gene evidence remains diffuse across many unrelated genes.
\item High retrieval confidence with low specificity confidence, which indicates a stable morphology neighborhood that still does not concentrate on target-like routed gene evidence.
\end{itemize}

In practice, the implementation can use specificity confidence as a guardrail rather than a passive annotation. In particular, opposite-polarity programs remain experimental and may be suppressed by default unless routed gene evidence reaches at least a moderate specificity level.

\section{Recommended defaults}
A sensible initial default is:
\begin{itemize}[leftmargin=2em]
\item reference bundle: JUMP pilot, matched cell line and timepoint,
\item query aggregation: median within user-supplied group,
\item similarity metric: cosine,
\item polarity: similar only,
\item neighbor restriction: top 20 hubness-adjusted reference neighbors with nonnegative similarity retained by default,
\item same-modality-first retrieval when query modality metadata is available,
\item adaptive neighborhood truncation after a minimum coherent local panel,
\item hubness penalty: rank-based downweighting when the bundle provides hub scores,
\item gene recurrence penalty: weak IDF-style downweighting of genes that recur across many reference perturbations,
\item modality balance: equal across available modalities,
\item selection: top-k with $K=200$,
\item normalization: within-set $\ell_1$,
\item GMT emission with 100--500 size bounds.
\end{itemize}

Opposite-polarity programs remain useful as a secondary exploratory view, but should be treated as experimental unless retrieval confidence and context matching are strong.

\begin{thebibliography}{9}

\bibitem{jump_main}
JUMP-Cell Painting Consortium.
JUMP-Cell Painting Consortium public portal.
\url{https://jump-cellpainting.broadinstitute.org/}

\bibitem{jump_gallery}
Cell Painting Gallery documentation: complete datasets.
\url{https://broadinstitute.github.io/cellpainting-gallery/complete_datasets.html}

\bibitem{jump_target}
JUMP-Target metadata repository.
\url{https://github.com/jump-cellpainting/JUMP-Target}

\bibitem{jump_pilot}
Chandrasekaran et al.
Three million images and morphological profiles of cells treated with matched chemical and genetic perturbations.
Nature Methods (2024).

\bibitem{jump_morphmap}
Chandrasekaran et al.
Morphological map of under- and overexpression of genes in human cells.
Nature Methods (2025).

\end{thebibliography}

\end{document}
